\section{Measurement protocol}\label{sec:protocol}

The claims in this paper are of a kind that is easy to state and hard to check, so
the apparatus is part of the contribution and is described before the results.
Three mechanisms do the work: unit-level provenance, a separation between claims
that must be true and claims that may be false, and an explicit refusal to
substitute estimates for measurements. The full recipe is in
\ref{app:artifact}.

\subsection{Unit-level provenance and crash-safe resume}

Each experiment is a \emph{unit} that writes exactly one JSON result file
containing the measurements, the configuration, the device, and two dictionaries:
\texttt{assertions} and \texttt{hypotheses}. A unit is marked \texttt{running}
before it starts and \texttt{done} only after its file exists and hashes to the
recorded digest; a killed process leaves \texttt{running}, which becomes
\texttt{interrupted} on the next launch and is re-executed. Completed units are
skipped by content hash. Every figure and table in this paper is generated from
these files at build time: the LaTeX source contains no typed measurement, and
the number of macros emitted by the build script equals the number of quantities
the text uses. This is what allows the sentence ``\VramPswaBig\,MiB at
$n=\LargestN$'' to be checkable rather than merely plausible.

A dependency fingerprint guards the derived layer. The figure stage records the
digests of the result files it consumed; if any upstream unit is re-run, the
digest changes, the figure unit is invalidated and regenerated. The mechanism was
added after it failed: a device benchmark was re-measured, the figure stage
reported \texttt{skipped} because its own output existed, and the repository
briefly contained a figure that did not correspond to the results beside it. A
protocol that can produce that failure silently is not adequate for a paper whose
argument is about measurement.

\subsection{Assertions versus hypotheses}\label{sec:audit}

The distinction is small and consequential. An \emph{assertion} is an internal
consistency requirement: that the mask is sparse, that the outputs are finite,
that two variants differ, that a bound holds. If an assertion is false the
artefact is broken and the run must fail. A \emph{hypothesis} is a scientific
claim: that sparsity saves memory, that a periodic anchor is supported by the
data, that adaptation helps. A hypothesis may be false, and when it is, the
correct outcome is a recorded refutation and a rewritten paper, not a failed run.

Conflating the two is not a cosmetic error. In an earlier version of this
apparatus both dictionaries were merged, so a refuted hypothesis made the
verification gate fail, and the cheapest way to a green run was to weaken the
claim until it passed. Separating them removed that incentive: the gate now
requires every assertion to be true and merely \emph{reports} the hypotheses that
came out false. Section~\ref{sec:attn_withdraw} lists what this caught in our own
work.

Two further rules apply to every unit. A measurement that could not be taken is
recorded as \texttt{blocked} with the command that would unblock it, and never
replaced by an estimate or a value borrowed from the literature. A quantity that
the pipeline could not compute, whether a parse failure, a missing class or a
degenerate split, is recorded as a failure and counted, never imputed;
Section~\ref{sec:lang} shows what imputation would have done to the headline
number of the language-model experiment.

\subsection{Device benchmark protocol}\label{sec:device_protocol}

Six decisions determine whether a memory measurement means anything, and we state
all six because three of them were wrong in our first attempt.

\begin{enumerate}\itemsep2pt
\item \emph{Baseline decontamination.} Peak memory is reported as the difference
between the maximum allocated bytes during the operation and the allocation
immediately before it, after \texttt{gc.collect()} and
\texttt{torch.cuda.empty\_cache()}, with the peak counter reset \emph{after} the
query, key and value tensors exist. Without the reset, a stale tensor inflated the
baseline by $8.5$\,MiB in our first run and produced impossible numbers at the
smallest length.
\item \emph{Two dense baselines.} We report a naive dense implementation, which
materialises the $n\times n$ score matrix, and a fused kernel
\cite{dao2022flashattention} that does not. The first is the baseline that
inflates a sparse saving; the second is the competitor a real system faces. Both
are reported for every length.
\item \emph{Two sparse implementations.} A batched variant (one gather plus one
batched matrix product per sequence, the way a production kernel would run) and a
block-streaming variant (minimal live memory, one launch per query block). The
second exists so that a memory claim cannot be dismissed as an artefact of an
inefficient gather. The two are reported separately and they behave very
differently, which is itself a result.
\item \emph{Omission is a datum.} An out-of-memory outcome is stored as
\texttt{status="oom"}, never skipped and never interpolated; the largest length at
which the dense path fails while a sparse path succeeds is reported as a measured
quantity.
\item \emph{Tail latency.} We report the median and the 95th percentile of
\GpuReps{} repetitions after warm-up, with a device synchronisation around each
repetition, because a mean over unsynchronised asynchronous launches is not a
latency.
\item \emph{Disclosure of implementation dependence.} The batched variant pads
every query block to the largest number of key blocks required by \emph{any} block
in the sequence, which is the last one. That choice is correct but not optimal,
and it inflates its memory relative to a hand-tuned kernel; we quantify the
consequence and scope the claim accordingly (Section~\ref{sec:attn_disclose}
rather than in a footnote).
\end{enumerate}

\subsection{Elicitation protocol for language models}

Two interfaces are measured on the same rows.
\emph{Logit interface}: the model is asked a yes/no question, and the score is
$p=P(\texttt{Yes})/\big(P(\texttt{Yes})+P(\texttt{No})\big)$ read at the answer
position of a single forward pass, so no text is parsed and no failure mode
exists. \emph{JSON interface}: the model is asked for a small JSON object, the
output is parsed with a regular expression, and every unparsable answer is counted
as a failure. The second interface is what practitioners do; the first exists to
separate the cost of the model from the cost of the interface.

Only the logits of the final position are computed, because materialising the full
vocabulary distribution for a batch allocates $\text{batch}\times\text{seq}\times
\lvert V\rvert$ floats and exhausts the device at the batch size we use. The
adaptation experiment replaces the hard-coded targets of the original pipeline
(which used the constants $0.90$ and $0.05$) with the per-row continuous
probability predicted by the reference model on the training split, so the target
has \TargetDistinct{} distinct values with standard deviation \TargetSd.

\subsection{Statistical protocol}

Systems are always compared on identical rows, and the comparison statistics are
paired. For a ranking metric we use a paired permutation test over rows, which
makes no distributional assumption about AUC-PR and is exact under exchangeability
\cite{efron1993bootstrap}; we report the permutation $p$-value and the observed
difference. For confidence intervals we bootstrap the difference with two
resampling units: rows, and entities (cards), because several transactions share a
card and the row unit understates the variance. The entity-level interval is the
statistically correct one and we report both so the reader can see the difference.
For paired decisions we use the exact McNemar test \cite{mcnemar1947note}, for
score differences the DeLong test \cite{delong1988comparing}, and we control the
family-wise error with Holm's step-down procedure \cite{holm1979simple}. The number
of training seeds is set by an a-priori power calculation rather than by
convention; the calculation is reported with its inputs in
\ref{app:tables}.

Calibration is reported with the Brier score and the expected calibration error
before and after Platt scaling fitted on validation only
\cite{guo2017calibration}. Fairness is reported as equalised-odds gaps
\cite{hardt2016equality} on the attributes the second dataset provides. Both are
reported even where they are unfavourable: Section~\ref{sec:system} reports a
calibration intervention that does not help and a drift comparison that is
confounded by prevalence.
